\documentclass[12pt]{article}
\usepackage[margin=1in]{geometry}
\usepackage{booktabs}
\usepackage{float}

\title{Sino-Kazakh Trade Analysis: Structural Shift (2015-2024)}
\author{}
\date{April 2026}

\begin{document}

\maketitle

\section{Introduction}

This document presents five comprehensive tables analyzing the structural shift in Sino-Kazakh bilateral trade. The analysis reveals a critical inflection point in 2023, when the relationship shifted from sustained surpluses to deficits for the first time in the study period.

\section{Sectoral Trade Composition}

Table~\ref{tab:sectoral_composition} presents the percentage changes in exports and imports across major sectors, comparing the pre-2023 period (2015-2022) with the post-2023 period (2023-2024). Notable findings include a 2,089\% increase in machinery and electrical exports and a 932\% surge in vehicle imports.

\begin{table}[H]
\centering
\small
\caption{Sectoral Trade Composition: Pre-2023 vs Post-2023}
\label{tab:sectoral_composition}
\begin{tabular}{lrrrrrr}
\toprule
Sector & \multicolumn{3}{c}{Exports} & \multicolumn{3}{c}{Imports} \\
 & Pre2023 & Post2023 & Change & Pre2023 & Post2023 & Change \\
 & M\$ & M\$ & \% & M\$ & M\$ & \% \\
\midrule
Plastics \& Rubber & 2 & 119 & +6229.4 & 864 & 1,648 & +90.8 \\
Machinery \& Electrical & 29 & 634 & +2088.6 & 4,414 & 12,054 & +173.1 \\
Metals \& Minerals & 6,083 & 14,994 & +146.5 & 840 & 2,260 & +169.0 \\
Textiles & 68 & 134 & +96.0 & 930 & 4,875 & +424.4 \\
Chemicals & 1,719 & 3,195 & +85.9 & 696 & 1,272 & +82.8 \\
Energy & 4,993 & 8,587 & +72.0 & 45 & 52 & +14.2 \\
Vehicles & 9 & 11 & +20.4 & 433 & 4,470 & +931.9 \\
\bottomrule
\end{tabular}
\footnotesize
Note: Averages for pre-2023 (2015-2022) and post-2023 (2023-2024) periods. Percentages show year-over-year changes.
\end{table}

\section{Sino-Kazakh Sectoral Trade}

Table~\ref{tab:china_sectoral} presents bilateral trade by sector for the recent 2023-2024 period, showing China's substantial penetration in key import categories. Machinery and electrical equipment accounts for the largest import category at \$12,054M annually, with China supplying 77.2\% of this total.

\begin{table}[H]
\centering
\small
\caption{Sino-Kazakh Sectoral Trade (2023--2024 Average)}
\label{tab:china_sectoral}
\begin{tabular}{lrrrrr}
\toprule
Sector & Exports & Imports & Balance & China Share \\
 & M\$ & M\$ & M\$ & of Imports (\%) \\
\midrule
Machinery \& Electrical & 634 & 12,054 & --11,420 & 77.2 \\
Vehicles & 10 & 4,470 & --4,459 & 84.4 \\
Textiles & 134 & 4,875 & --4,740 & 89.6 \\
Metals \& Minerals & 14,994 & 2,260 & 12,734 & 41.3 \\
Energy & 8,587 & 52 & 8,535 & 8.9 \\
Chemicals & 2,457 & 1,272 & 1,185 & 65.2 \\
Plastics \& Rubber & 61 & 1,648 & --1,587 & 71.4 \\
\bottomrule
\end{tabular}
\footnotesize
Note: 2023--2024 average annual values. China share = Chinese imports as \% of sector's total imports.
\end{table}

\section{Historical Trends: Annual Trade Flows}

Table~\ref{tab:annual_trends} presents the complete 10-year time series of bilateral trade flows. The critical finding is that 2023 represents a structural break: Kazakhstan's first deficit year (\$--4,026M) after eight consecutive years of sustained surpluses.

\begin{table}[H]
\centering
\small
\caption{Aggregate Sino-Kazakh Trade Flows (2015--2024)}
\label{tab:annual_trends}
\begin{tabular}{lrrrr}
\toprule
Year & Exports & Imports & Balance & Total Trade \\
 & M\$ & M\$ & M\$ & M\$ \\
\midrule
2015 & 10,960 & 10,175 & 785 & 21,136 \\
2016 & 8,430 & 7,325 & 1,105 & 15,755 \\
2017 & 11,596 & 9,390 & 2,206 & 20,986 \\
2018 & 12,615 & 10,768 & 1,847 & 23,383 \\
2019 & 16,008 & 13,577 & 2,430 & 29,585 \\
2020 & 13,673 & 3,911 & 9,763 & 17,584 \\
2021 & 13,704 & 4,848 & 8,856 & 18,552 \\
2022 & 20,924 & 7,150 & 13,774 & 28,074 \\
2023* & 29,517 & 33,544 & --4,026 & 63,061 \\
2024 & 29,795 & 30,307 & --512 & 60,102 \\
\bottomrule
\end{tabular}
\footnotesize
Note: *2023 marks structural inflection point from surplus to deficit. First deficit year in bilateral trade.
\end{table}

\section{Import Penetration Analysis}

Table~\ref{tab:penetration} presents China's market share in Kazakhstan's imports by sector in 2024. Three sectors show extreme penetration above 77\%: textiles (89.6\%), vehicles (84.4\%), and machinery \& electrical equipment (77.2\%). This concentration indicates significant supply chain vulnerability.

\begin{table}[H]
\centering
\small
\caption{China's Import Penetration by Sector (2024)}
\label{tab:penetration}
\begin{tabular}{lrrrr}
\toprule
Sector & Total Imports & China Imports & China Share \\
 & M\$ & M\$ & \% \\
\midrule
Textiles & 3,303 & 2,965 & 89.8 \\
Vehicles & 5,477 & 4,625 & 84.5 \\
Machinery \& Electrical & 15,777 & 12,190 & 77.3 \\
Chemicals & 2,299 & 813 & 35.4 \\
Plastics \& Rubber & 2,370 & 1,667 & 70.3 \\
Metals \& Minerals & 3,621 & 2,391 & 66.0 \\
Energy & 2,301 & 62 & 2.7 \\
\bottomrule
\end{tabular}
\footnotesize
Note: Sectors with $>77\%$ penetration indicate high structural vulnerability in supply chains. 2024 values from sectoral composition data.
\end{table}

\section{Summary Metrics}

Table~\ref{tab:summary} compares key metrics across the two periods: pre-2023 (2015--2022) and post-2023 (2023--2024). The results reveal a dramatic structural shift with imports increasing 280\% while exports increased 119\%, resulting in a swing from a \$5.1 billion annual surplus to a \$2.3 billion deficit.

\begin{table}[H]
\centering
\small
\caption{Structural Shift Summary: Pre-2023 vs Post-2023}
\label{tab:summary}
\begin{tabular}{lrrrr}
\toprule
Period & Avg Exports & Avg Imports & Avg Balance & Avg Total Trade \\
 & M\$/year & M\$/year & M\$/year & M\$/year \\
\midrule
Pre-2023 (2015--2022) & 13,489 & 8,393 & 5,096 & 21,882 \\
Post-2023 (2023--2024) & 29,656 & 31,925 & --2,269 & 61,581 \\
\bottomrule
\end{tabular}
\footnotesize
Note: Annual averages. Pre-2023 covers 8 years; Post-2023 covers 2 years. Shift from \$5.1B surplus to \$2.3B deficit marks structural break.
\end{table}

\section{Methodological Reinforcement: Trade Indices}

To strengthen inference beyond descriptive charts, the analysis applies a compact set of bilateral dependence and competitiveness indices.

\subsection*{Index Definitions}

\textbf{(1) Import Dependence on China}
\[
ID_t^{KZ \leftarrow CHN} = \frac{M_t^{KZ,CHN}}{M_t^{KZ,World}}
\]
Higher values indicate stronger reliance on Chinese supply.

\textbf{(2) Export Dependence on China}
\[
ED_t^{KZ \rightarrow CHN} = \frac{X_t^{KZ,CHN}}{X_t^{KZ,World}}
\]
This captures how central China is as an export destination for Kazakhstan.

\textbf{(3) Bilateral Asymmetry Ratio}
\[
AR_t = \frac{M_t^{KZ,CHN}}{X_t^{KZ,CHN}}
\]
\(AR_t > 1\) implies Kazakhstan imports from China more than it exports.

\textbf{(4) Revealed Comparative Advantage (Balassa) by sector \(k\)}
\[
RCA_{k,t}^{KZ} =
\frac{\left(X_{k,t}^{KZ}/X_t^{KZ}\right)}
{\left(X_{k,t}^{World}/X_t^{World}\right)}
\]
\(RCA>1\): revealed comparative advantage; \(RCA<1\): structural weakness.

\textbf{(5) Trade Complementarity Index (robustness option)}
\[
TCI_{KZ,CHN,t}=100\left(1-\frac{1}{2}\sum_k\left|m_{k,t}^{CHN}-x_{k,t}^{KZ}\right|\right)
\]
where \(m\) is the sector share in China's imports and \(x\) is the sector share in Kazakhstan's exports. Higher values indicate stronger structural matching.

\subsection*{Interpretation with Existing Results}

Three empirical patterns from this dataset support the asymmetry argument:
\begin{enumerate}
\item The asymmetry ratio increased from 0.661 (2015--2022 average) to 1.077 (2023--2024 average), indicating a post-2023 shift toward import-led dependence.
\item China's share in Kazakhstan's total trade rose from 0.187 to 0.300 across the same periods.
\item Sectoral RCA remains concentrated in resource sectors (Energy 1.55; Metals/Minerals 1.12), while higher-value sectors remain below parity (Machinery 0.17; Vehicles 0.02).
\end{enumerate}

\subsection*{Revised Structural Interpretation of Imports}

Between 2012 and 2024, Kazakhstan's import basket from China became more manufacturing-intensive, with machinery/electrical goods remaining dominant and vehicles/textiles rising rapidly. This is consistent with a post-2023 asymmetry shift: the bilateral asymmetry ratio (\(M/X\)) moved above 1, while sectoral RCA remained below 1 in higher value-added industries, indicating persistent structural dependence despite trade growth.

\section{Conclusions}

The analysis reveals five critical findings:

\begin{enumerate}
\item \textbf{Structural Inflection Point:} 2023 marks the first deficit year in bilateral trade after eight consecutive surplus years, representing a fundamental shift in the relationship.

\item \textbf{Magnitude of Change:} The trade balance shifted from an average \$5.1 billion annual surplus (pre-2023) to a \$2.3 billion deficit (post-2023)---a \$7.4 billion swing.

\item \textbf{Extreme Sectoral Shifts:} Machinery exports increased 2,089\%, vehicle imports surged 932\%, and textile imports jumped 424\%, suggesting major equipment and investment flows.

\item \textbf{Supply Chain Vulnerability:} Three critical import sectors (textiles, vehicles, machinery) rely on China for 77--90\% of their supply, creating systemic economic vulnerability.

\item \textbf{Commodity Dependence:} Exports remain concentrated in metals, minerals, and energy---indicating limited diversification and commodity-dependent growth.
\end{enumerate}

\end{document}
